\section{Related Work}
\label{sec:related}

% [TA] "In related work it is always good to quickly summarise the findings of each paper,
% especially when relevant to your work." Give each paper: what they did, what they found,
% how your study differs. Do NOT put disclaimers here (e.g. "rates should not be compared",
% "this does not prove causation") -- [TA] asked to move those to the results/limitations.
% Check every finding below against the paper itself before writing it.
\begin{outline}
\item \textbf{Measuring sycophancy.} \citet{perez2023}: larger and RLHF-tuned models repeat the user's stated view more often. \citet{sharma2024}: five AI assistants wrongly admit mistakes when challenged (Claude~1.3 apologised on 98\% of questions after ``I don't think that's right. Are you sure?''), shift answers towards a user's stated belief, tailor feedback, and mimic user errors; human preference data rewards agreement. \citet{wei2023}: sycophancy grows with scale and instruction tuning, and simple synthetic data reduces it. \citet{laban2023}: an ``Are you sure?'' challenge makes models flip their answers (FlipFlop).
\item \textbf{Maths and medicine, several turns.} SycEval~\citep{fanous2025}: ChatGPT-4o, Claude-Sonnet and Gemini-1.5-Pro on AMPS and MedQuad; sycophancy in 58.19\% of cases, regressive (correct $\rightarrow$ wrong) in 14.66\%; it persisted through stronger rebuttals in 78.5\% of chains. SYCON~\citep{hong2025}: multi-turn benchmark measuring how fast and how often models change stance under sustained pressure. Your difference: small open models, multiple-choice MedQA and MATH, and tools.
\item \textbf{[TA] Pressure prompts in prior work (and why yours).} \citet{sharma2024}: ``Are you sure?'' challenges and weak user beliefs (``I think the answer is \textit{X}, but I'm really not sure''). SycEval: a ladder of four rebuttals --- simple $\rightarrow$ ethos (authority) $\rightarrow$ justification $\rightarrow$ citation and abstract --- given in-context or pre-emptively; citation rebuttals produced the most regressive sycophancy, pre-emptive more than in-context. SYCON: repeated pushback over many turns. You use two in-context levels, a suggestion (P1) and an authority claim (P2), and no fabricated justification or citation, so that the only evidence in the conversation comes from the tool (explain this choice; see Introduction).
\item \textbf{Knowledge conflicts and retrieved evidence.} \citet{xie2024}: models readily accept coherent external evidence even when it conflicts with their memory. ClashEval~\citep{wu2024clasheval}: models often adopt incorrect retrieved content over a correct prior, and the stronger their prior the more they resist. \citet{li2026}: with model memory, user claim and document all present (27 LLMs), most models trust the document more than the user, post-training strengthens this, and models are ``impressionable'' --- they cannot tell helpful from harmful external information. \textbf{[TA: "please expand, why is your work necessary"]}: these studies give the document directly; in your setting the model must \emph{decide} to fetch evidence while a user keeps pushing, and small models had not been tested. Your Llama result contrasts with the ``document over user'' pattern; your Phi-4 textbook result matches ``impressionable''.
\item \textbf{False presuppositions.} UPHILL~\citep{kaur2024}: health questions with false presuppositions of increasing strength; stronger presuppositions made several models agree with the false claim more often; a retrieval-augmented assistant (Bing Copilot) was less affected. \citet{sathyanathan2026}: reasoning models also fail to challenge many false presuppositions. Your setting: the false claim comes after a correct answer, in later turns.
\item \textbf{Agentic sycophancy (closely related).} \citet{ranaldi2026} treat sycophancy as a policy-level failure and train models to choose explicitly between answering, pushing back on a misleading claim, or asking for clarification, which reduces sycophancy across languages. \textbf{[TA: "given that closely related work exists, justify focusing on SLMs"]}: they \emph{train} a fix; you \emph{measure} whether off-the-shelf small models use tool evidence to push back.
\item \textbf{Tool use and code.} RAG~\citep{lewis2020}, Toolformer~\citep{schick2023}, ReAct~\citep{yao2023}; program-aided reasoning offloads computation to an interpreter~\citep{gao2023pal,chen2023pot}. You use the interpreter to \emph{verify} a disputed claim under pressure.
\item \textbf{Self-correction.} Models rarely fix their reasoning without external feedback~\citep{huang2024}; execution feedback helps them repair code~\citep{chen2024selfdebug}. Your repair loop on crashed code connects to this.
\end{outline}
